\chapter*{Resumen}
\addcontentsline{toc}{chapter}{Resumen}
\markboth{Resumen}{Resumen}
Este trabajo estudia la capacidad de distintos algoritmos de aprendizaje automático para predecir la evolución de activos financieros en mercados con comportamientos muy variables. Se analizan modelos clásicos y modelos actuales, evaluando su rendimiento sobre distintos horizontes temporales y en activos como Bitcoin, Tesla y Nvidia. El proyecto incorpora señales de sentimiento específicas por activo a partir de noticias: para los tres activos se usa el tono medio de GDELT (Average Tone), junto con su intensidad de cobertura como proxy de volumen. El objetivo es comprobar si la combinación de patrones históricos con información de sentimiento permite mejorar la precisión de las predicciones. La evaluación se realiza con validación walk-forward, métricas RMSE/MAE/MAPE y métricas direccionales (acierto, F1 y MCC), junto con un test de Wilcoxon para comparar la ablación con y sin sentimiento. Los resultados muestran que los baselines simples siguen siendo difíciles de batir: ARIMA obtiene el mejor MAPE medio y Naive queda muy cerca; por activo, Naive lidera en BTC-USD y ARIMA en TSLA y NVDA. Los modelos de árboles (Random Forest y XGBoost) no mejoran a los baselines y, en esta configuración experimental, las arquitecturas de aprendizaje profundo (CNN, LSTM y GRU) presentan errores claramente superiores. En el experimento de horizontes sobre Bitcoin, todos los modelos empeoran al alargar el horizonte y GRU es el más consistente dentro de los modelos DL. El impacto del sentimiento es pequeño y heterogéneo: reduce el error en Tesla y en Nvidia con LSTM, pero lo incrementa en Bitcoin y en Nvidia con CNN/GRU; en ningún caso aparecen diferencias estadísticamente significativas en MAPE según el test de Wilcoxon. El test de Diebold-Mariano no encuentra mejoras significativas frente a Naive y, en varias comparaciones, sí detecta empeoramientos significativos, reforzando la necesidad de baselines y tests estadísticos antes de asumir ventaja predictiva.

\vspace{0.5cm}

\textbf{Palabras clave:} predicción financiera, aprendizaje automático, series temporales, sentimiento de mercado.

\chapter*{Abstract}
\addcontentsline{toc}{chapter}{Abstract}
\markboth{Abstract}{Abstract}
This work examines the performance of several machine learning algorithms applied to the prediction of financial assets in markets with highly variable behaviour. Both traditional and modern models are evaluated across different forecasting horizons, using assets such as Bitcoin, Tesla and Nvidia as case studies. The project introduces asset-specific sentiment signals from news: for all three assets it uses GDELT Average Tone, plus coverage intensity as a volume proxy. The goal is to assess whether combining historical patterns with external sentiment information improves prediction accuracy. Evaluation uses walk-forward validation, RMSE/MAE/MAPE, and directional metrics (hit rate, F1 and MCC), plus a Wilcoxon test for the with/without-sentiment ablation. Results show that simple baselines remain hard to beat: ARIMA achieves the lowest average MAPE and Naive is very close; by asset, Naive leads on BTC-USD while ARIMA leads on TSLA and NVDA. Tree-based models (Random Forest and XGBoost) do not outperform the baselines and, under this experimental configuration, deep learning architectures (CNN, LSTM and GRU) yield clearly higher errors. In the horizon experiment on Bitcoin, all models degrade as the horizon increases and GRU is the most consistent within the DL family. The sentiment impact is small and heterogeneous: it reduces error for Tesla and for Nvidia with LSTM, but increases error for Bitcoin and for Nvidia with CNN/GRU; Wilcoxon tests do not find statistically significant MAPE differences. The Diebold--Mariano test does not find significant improvements over Naive and, in several comparisons, detects statistically significant underperformance, reinforcing the need for baselines and statistical testing before claiming predictive advantage.

\vspace{0.5cm}

\textbf{Keywords:} financial prediction, machine learning, time series, market sentiment.
